\section{Modelos de lenguaje: de los vectores de palabras al posentrenamiento}

Las dos secciones anteriores explicaron los paradigmas de modelos y la base de entrenamiento; ahora volvemos a la puerta de entrada que hoy resulta más familiar para la mayoría: los modelos de lenguaje. Esta sección une en una sola línea «representación, preentrenamiento y posentrenamiento». La evolución de los modelos de lenguaje puede verse en tres etapas: primero, convertir las palabras en vectores; después, aprender las regularidades del lenguaje mediante el preentrenamiento; y, por último, usar instrucciones, preferencias y posentrenamiento de código abierto para que el modelo sea utilizable, controlable y capaz de dar servicio a los usuarios.

\begin{figure}[H]
\centering
\includegraphics[width=0.96\textwidth]{language-model-lineage.png}
\caption{Genealogía de los modelos de lenguaje: de los vectores de palabras al preentrenamiento y, después, a la alineación por instrucciones y al posentrenamiento de código abierto. Diagrama conceptual propio, elaborado a partir de la conversación de 02:21:29--03:08:08.}
\end{figure}

\begin{knowledgebox}{Lectura de la figura: el eje de los modelos de lenguaje va «de la representación a la conducta»}
Word2Vec resuelve la representación de las palabras; la traducción automática neuronal demostró que los modelos profundos pueden dar servicio en línea; GPT/BERT establecieron el paradigma del preentrenamiento; GPT-3 mostró la emergencia de capacidades con la escala; e InstructGPT y Tulu 3 llevaron los modelos hacia un uso práctico para las personas y hacia el posentrenamiento de código abierto.
\end{knowledgebox}

\subsection{Word2Vec y Google Translate: representación y despliegue}

Esta sección comienza con los antecedentes de los modelos de lenguaje, porque sin el aprendizaje de representaciones no habría existido el aprendizaje en contexto posterior. Word2Vec usa aprendizaje automático para convertir palabras en vectores, de modo que las relaciones semánticas entran en un espacio continuo. El despliegue de redes neuronales en Google Translate mostró, a su vez, que el aprendizaje profundo no solo funciona en los artículos académicos, sino que también puede sustituir a los métodos tradicionales en sistemas en línea a gran escala. Estos dos pasos respondieron, respectivamente, a «cómo representar el lenguaje» y a «cómo un modelo neuronal puede dar servicio a usuarios reales».

\begin{importantbox}{El significado de los vectores de palabras}
Los vectores de palabras proyectan palabras discretas en un espacio continuo, de modo que las palabras similares quedan más cerca en el espacio vectorial. Esto preparó el camino para las representaciones contextuales, el embedding del Transformer y la búsqueda semántica posteriores; sin embargo, sigue siendo una representación estática que no puede entender el significado de una palabra de forma dinámica según el contexto, como lo hace un LLM.
\end{importantbox}

\subsection{GPT-1, BERT, GPT-2 y GPT-3}

A partir de los vectores de palabras y la traducción en línea, los modelos de lenguaje entran de lleno en la era del preentrenamiento. GPT-1, BERT, GPT-2 y GPT-3, que se tratan en esta subsección, representan etapas distintas: GPT-1 propuso un nuevo paradigma de NLP basado en preentrenamiento no supervisado más fine-tuning supervisado; BERT, con masked language modeling bidireccional, se convirtió en «el rey de su época» en las tareas de comprensión; GPT-2 reforzó la ruta del preentrenamiento generativo e hizo que se empezara a considerar en serio la posibilidad de «prescindir del fine-tuning»; y GPT-3 combinó la Scaling Law, la apuesta organizacional y la ingeniería a gran escala, hasta convertirse en la señal más importante en la víspera de ChatGPT.

\begin{knowledgebox}{La bifurcación entre GPT y BERT}
\begin{tabularx}{\textwidth}{p{0.22\textwidth}p{0.34\textwidth}X}
\toprule
Ruta & Método de entrenamiento & Ventaja típica \\
\midrule
BERT & Predicción de palabras enmascaradas, contexto bidireccional & Destaca en tareas discriminativas como comprensión, clasificación y extracción. \\
GPT & next-token prediction autorregresiva & Gran capacidad de extensión a generación, completado, diálogo y tareas generales. \\
GPT-3 & Preentrenamiento autorregresivo a gran escala & El aprendizaje con pocos ejemplos, el aprendizaje en contexto y las capacidades emergentes pasan a ser el centro de atención. \\
\bottomrule
\end{tabularx}
\end{knowledgebox}

\subsection{InstructGPT y Tulu 3: llevar el modelo a la sociedad civilizada}

El significado central de InstructGPT es haber llevado al modelo de «saber continuar texto de internet» a «poder responder según la intención humana». Detrás de esto están métodos de posentrenamiento como los datos de instrucciones, los datos de preferencias y el RLHF. Tulu 3, por su parte, representa el intento de la comunidad de código abierto de sistematizar y hacer transparente el proceso de posentrenamiento, para que las capacidades de los modelos no queden solo en manos de unos pocos laboratorios de código cerrado.

\begin{warningbox}{El posentrenamiento no es un simple complemento}
El preentrenamiento le da al modelo conocimiento del mundo y capacidad lingüística; el posentrenamiento determina cómo responde a las personas, cómo sigue instrucciones, cómo rechaza solicitudes peligrosas, cómo produce salidas con el formato indicado y cómo completa tareas. Sin posentrenamiento, incluso un modelo potente puede tener dificultades para convertirse en un buen producto.
\end{warningbox}

\subsection{Resumen de la sección}

La línea de los modelos de lenguaje llevó a la IA del aprendizaje de representaciones a una puerta de entrada de interacción general. De Word2Vec a GPT-3 hay un crecimiento de capacidades; de InstructGPT a Tulu 3, un crecimiento en facilidad de uso, alineación y posentrenamiento de código abierto. Los productos de LLM actuales se construyen en el punto donde se cruzan estas dos curvas de crecimiento.
